\section{Does it hold for the released models?}
\label{sec:published}

\paragraph{The DNABERT family is ordered on three genomes.}
Table~\ref{tab:dnabert} fine-tunes the four released $k$-mer checkpoints on three real datasets with
three seeds each. Both the smooth and the thresholded statistic rise monotonically with $k$.

\begin{table}[t]
\centering
\caption{Released checkpoints, next-token fine-tuning, round~50, canary planted $16\times$, three
seeds. Canary surprise in \bits{} and verbatim extraction rate.}
\label{tab:dnabert}
\begin{tabular}{lcccccc}
\toprule
& \multicolumn{2}{c}{\emph{E.~coli}} & \multicolumn{2}{c}{yeast} & \multicolumn{2}{c}{GUE promoters} \\
\cmidrule(lr){2-3}\cmidrule(lr){4-5}\cmidrule(lr){6-7}
model & \bits & exact & \bits & exact & \bits & exact \\
\midrule
DNABERT \kmer{3} & 1.60 & 0\%   & 0.76 & 68\%  & 1.53 & 5\%   \\
DNABERT \kmer{4} & 0.88 & 33\%  & 0.56 & 78\%  & 1.00 & 23\%  \\
DNABERT \kmer{5} & 0.59 & 55\%  & 0.36 & 100\% & 0.51 & 68\%  \\
DNABERT \kmer{6} & 0.30 & 98\%  & 0.21 & 100\% & 0.29 & 90\%  \\
\midrule
DNABERT-2 (BPE)  & 0.05 & 100\% & 0.05 & 100\% & 0.04 & 100\% \\
HyenaDNA medium  & 0.12 & 77\%  & 0.12 & 82\%  & 0.11 & 83\%  \\
\bottomrule
\end{tabular}
\end{table}

\paragraph{Two models break the ordering from outside the tokeniser axis.}
DNABERT-2 memorises more than any $k$-mer checkpoint, and HyenaDNA, which is character level,
memorises more than the \kmer{5}. Neither contradicts Sec.~\ref{sec:controlled}. DNABERT-2 has 117M
parameters against 86--89M, ALiBi instead of learned absolute positions, and a different pretraining
corpus; on a matched backbone a BPE vocabulary of the same average piece length sits between the
\kmer{3} and the \kmer{6}. HyenaDNA differs in architecture and pretraining, not in tokenisation.
This is precisely why the causal claim needs the controlled grid: the released family alone cannot
support it, and reading these rows as a tokeniser effect would be the mistake we are trying to
correct.

\paragraph{On human chromosome~22 the $k$-mer family is not ordered, and we say so.}
The fourth dataset breaks the pattern of Table~\ref{tab:dnabert}, and four independent draws of the
canary split do not resolve it. At equal exposure (round~50) \excess{} rises weakly with $k$ ---
\NUM{1.45}, \NUM{1.51}, \NUM{1.50}, \NUM{1.68} for $k=3,4,5,6$, with an across-draw spread of
\NUM{0.04}--\NUM{0.07} that separates the \kmer{6} from the \kmer{3} but leaves $k=4$ and $k=5$
tied. At matched held-out loss the direction reverses: \NUM{1.48}, \NUM{1.06}, \NUM{1.01},
\NUM{1.37}, with the \kmer{3} highest, because it overfits chromosome~22 most slowly and matching
the loss gives it the most exposure to the planted records. The six rising branches overlap in a
window \NUM{0.005}~\bits{} wide, which leaves no ground to prefer either normalisation, and the
extraction rate is not monotone either (\NUM{94}, \NUM{87}, \NUM{78}, \NUM{88\%}). \textbf{We
therefore base no tokeniser ordering on chromosome~22.} What is robust there comes from outside the
tokeniser axis: DNABERT-2 (\NUM{2.35}) and HyenaDNA (\NUM{2.39}) sit above the entire $k$-mer family
under both readings and on every draw, as on the other three datasets.

\paragraph{Why the canary split is drawn four times.}
Which canaries are members is fixed by one data seed, so within an arm it is the same split for
every model, and with real haplotypes the pretrained checkpoints already read the member group as
easier before any fine-tuning --- up to \NUM{0.9}~\bits{} on a single draw. Re-drawing the split
four times shows that this offset must \emph{not} be subtracted: it is gone after one epoch
(\excess{} at round~1 is \NUM{0.00}$\pm$\NUM{0.05} on every draw), it propagates with a slope of
only \NUM{0.05}--\NUM{0.19}, and subtracting it inflates the across-draw spread from
\NUM{0.04}--\NUM{0.11} to \NUM{0.15}--\NUM{0.39} (Appendix~A9). Every number in this section and in
Sec.~\ref{sec:human} is therefore uncorrected, with the across-draw spread as its error bar.

\paragraph{Chromosome 21 says the same thing, including the part we cannot explain.}
We repeated the whole panel on a second chromosome, with canaries built from chromosome~21 by the
same filter (400 segments, 365 carriers) and two seeds for the $k$-mer family. Everything that
matters reproduces. Rare-allele recovery at $r{=}16$ is \NUM{0.84}--\NUM{1.00} against a
\NUM{0.23}--\NUM{0.29} control; DNABERT-2 again leaks a record seen \emph{once}, recovering
\NUM{77\%} of its rare alleles against a \NUM{27\%} control and reproducing \NUM{44\%} of those
segments verbatim, while no $k$-mer checkpoint extracts a single-copy record. And the anomaly
reproduces too: verbatim extraction at $r{=}16$ is \NUM{79}, \NUM{68}, \NUM{48}, \NUM{91\%} for
$k=3,4,5,6$, the same non-monotone shape with the same \kmer{5} dip we report on chromosome~22. An
artefact of one chromosome would not survive the change of chromosome; we still have no mechanism
for it, and we still base no ordering on it.
